\origpage{384}{N. Robertson, P.D. Seymour / Journal of Combinatorial Theory, Series B 90 (2004) 325--385}
\textbf{Proof of 4.1.} Suppose that some $\chi$ satisfying $\mathbf{S}_1\text{--}\mathbf{S}_4$ is orientedly bad. Then (by replacing the $\Omega_\chi(k)$'s and $\Omega_\chi(S)$'s by isomorphic well-quasi-orders) we may choose $\chi$ to be disjoint. Let $(\Gamma_0, \phi_0), (\Gamma_1, \phi_1), \dots$ be a similarly oriented bad sequence for $\chi$. By 6.1 we may assume that for all $j > i \ge 0$ there is no inflation of $(\Gamma_i, \phi_i)$ in $(\Gamma_j, \phi_j)$. Choose $n \ge 4$ and even such that 15.2 holds (for $(\Gamma_0, \phi_0)$). By 5.4 and 5.5, we may assume (by replacing the sequence $(\Gamma_i, \phi_i)$ ($i \ge 1$) by a subsequence) that $dist(\Gamma_i), rep(\Gamma_i) > 25n \cdot 5^{2|E(\Gamma_0)|}$ for all $i \ge 1$. From 15.2 we obtain

\hspace*{12pt}(1) \textit{For each $i \ge 1$, and for every tangle $\mathscr{T}$ in $\Gamma_i$, and every tie-breaker $\lambda$ in $\Gamma_i$, either}

\hspace*{12pt}$\bullet$ \textit{$rep(\mathscr{T}) \le 25n \cdot 5^{2|E(\Gamma_0)|}$, or}

\hspace*{12pt}$\bullet$ \textit{$(\Gamma, \phi, \lambda, \mathscr{T})$ is $25n \cdot 5^{2|E(\Gamma_0)|}$-flawed in distance, or}

\hspace*{12pt}$\bullet$ \textit{$(\Gamma, \phi, \lambda, \mathscr{T})$ is $\frac{1}{2}n(|E(\Gamma_0)| + 1)$-flawed in freedom, or}

\hspace*{12pt}$\bullet$ \textit{for some long side $S$ of $\Phi_\chi$, let $f_1, \dots, f_k$ be the edges of $\Gamma_0$ bordering $S$ in order; then $(\Gamma, \phi, \lambda, \mathscr{T})$ is $(\frac{1}{2}n, (\phi_0(f_1), \phi_0(f_2), \dots, \phi_0(f_k)))$-flawed on $S$, or}

\hspace*{12pt}$\bullet$ \textit{for some internal $e \in E(\Gamma_0)$, $(\Gamma, \phi, \lambda, \mathscr{T})$ is $(|E(\Gamma_0)|, n, \phi_0(e))$-flawed internally.}\fixnote{fix:4.1-claim1}{This claim also holds with the strict bounds of pages 357 and 362. In the last two outcomes, $rep(\mathscr{T}) > 25n \cdot 5^{2|E(\Gamma_0)|}$ since the first outcome fails, and $dist(\Gamma_i) > 25n \cdot 5^{2|E(\Gamma_0)|}$. For ``flawed internally'', with $m = |E(\Gamma_0)|$, we have $5^{2m+2}n = 25n \cdot 5^{2|E(\Gamma_0)|}$.}

Let $\mathscr{S}$ be a similarly oriented set of $\chi$-coloured paintings containing $(\Gamma_i, \phi_i)$ for each $i \ge 0$, and closed under $\Phi_\chi$-preserving homeomorphisms of $\Sigma_\chi$. This exists since $\{(\Gamma_i, \phi_i) \colon i \ge 0\}$ is similarly oriented and so are all images of each $(\Gamma_i, \phi_i)$ under $\Phi_\chi$-preserving homeomorphisms. Let $\mathscr{D}$ be the set of all $(\Gamma, \phi, \lambda, \mathscr{T})$ such that $(\Gamma, \phi) \in \mathscr{S}$, $\lambda$ is a tie-breaker in $\Gamma$, and $\mathscr{T}$ is a tangle in $\Gamma$.

By 10.3, 10.4, 10.5, 11.1 and 14.1, there is a well-behaved set $\mathscr{C}$ of $\chi$-places such that for each $(\Gamma, \phi, \lambda, \mathscr{T}) \in \mathscr{D}$, if $ord(\mathscr{T}) > 25n \cdot 5^{2|E(\Gamma_0)|}$ then there is a rooted location $\mathscr{L}$ which $(25n \cdot 5^{2|E(\Gamma_0)|} + 1)$-isolates $\mathscr{T}$ and for which $(\Gamma, \phi, \mathscr{L}) \in \mathscr{C}$. By 9.2 there exist $j > i \ge 1$ such that there is an inflation of $(\Gamma_i, \phi_i)$ in $(\Gamma_j, \phi_j)$. This is a contradiction, and completes the proof. $\square$

\sectionhead{References}

[1] N. Robertson, P.D. Seymour, Graph Minors. VI. Disjoint paths across a disc, J. Combin. Theory Ser. B 41 (1986) 115--138.

[2] N. Robertson, P.D. Seymour, Graph Minors. VII. Disjoint paths on a surface, J. Combin. Theory Ser. B 45 (1988) 212--254.

[3] N. Robertson, P.D. Seymour, Graph Minors IV. Tree-width and well-quasi-ordering, J. Combin. Theory. Ser. B 48 (1990) 227--254.

[4] N. Robertson, P.D. Seymour, Graph Minors. VIII. A Kuratowski theorem for general surfaces, J. Combin. Theory Ser. B 48 (1990) 255--288.

[5] N. Robertson, P.D. Seymour, Graph Minors. X. Obstructions to tree-decomposition, J. Combin. Theory Ser. B 52 (1991) 153--190.

[6] N. Robertson, P.D. Seymour, Graph Minors. XI. Circuits on a surface, J. Combin. Theory Ser. B 60 (1994) 72--106.

[7] N. Robertson, P.D. Seymour, Graph Minors. XII. Distance on a surface, J. Combin. Theory Ser. B 64 (1995) 240--272.
